\section*{Chapter \ref{ch:g10:heart-lungs-effort} --- Heart and Lungs During Effort}

\begin{solution}{exo:g10:heart-lungs-effort:1}
$0.6 \times 14 = \qty{8.4}{L/min}$; $2.2 \times 35 = \qty{77}{L/min}$.
\end{solution}

\begin{solution}{exo:g10:heart-lungs-effort:2}
The volume of blood pumped by one side of the heart per minute:
$Q = f \times V_s$. $65 \times 0.075 \approx \qty{4.9}{L/min}$.
\end{solution}

\begin{solution}{exo:g10:heart-lungs-effort:3}
Right atrium, right ventricle, pulmonary artery, lungs (pulmonary
circulation), pulmonary veins, left atrium, left ventricle, aorta,
organs (systemic circulation), veins, right atrium.
\end{solution}

\begin{solution}{exo:g10:heart-lungs-effort:4}
Tidal volume \qty{0.5}{L}; vital capacity \qty{4.5}{L}; residual
volume \qty{1.2}{L}.
\end{solution}

\begin{solution}{exo:g10:heart-lungs-effort:5}
About $220 - 16 = 204$ and $220 - 60 = 160$ beats per minute.
\end{solution}

\begin{solution}{exo:g10:heart-lungs-effort:6}
\qty{100}{W}: $120 \times 0.105 \approx \qty{12.6}{L/min}$;
\qty{300}{W}: $198 \times 0.115 \approx \qty{22.8}{L/min}$; a factor
of 1.8.
\end{solution}

\begin{solution}{exo:g10:heart-lungs-effort:7}
Rest: \qty{2.5}{L/min}; exercise: $0.04 \times 24 \approx
\qty{1.0}{L/min}$. The share falls twelvefold, the absolute flow only
by a factor of 2.6, because the output itself has grown.
\end{solution}

\begin{solution}{exo:g10:heart-lungs-effort:8}
$18 \times 0.14 = \qty{2.5}{L/min}$: about three quarters of the
\qty{3.45}{L/min} of a trained student --- close to it for a smaller
person.
\end{solution}

\begin{solution}{exo:g10:heart-lungs-effort:9}
The brain cannot tolerate even seconds without oxygen and glucose, so
its supply is protected. The gut and kidneys can suspend digestion and
filtration for an hour without harm, and their share is lent to the
muscles.
\end{solution}

\begin{solution}{exo:g10:heart-lungs-effort:10}
$V_s = 5000/42 \approx \qty{120}{mL}$, against about \qty{70}{mL}: her
enlarged heart delivers the resting output in far fewer beats.
\end{solution}

\begin{solution}{exo:g10:heart-lungs-effort:11}
The brain's command to move, and its anticipation of the effort, are
sent to the cardiac centres and the adrenal glands before the muscles
act. A runner relying only on carbon dioxide sensing would start with
a resting output, run the first minute on \omterm{def:g10:cell-metabolism:fermentation}{fermentation}, and be forced
to slow by lactic acid before the delivery caught up.
\end{solution}

\begin{solution}{exo:g10:heart-lungs-effort:12}
$195 \times 0.110 \times 0.150 \approx \qty{3.2}{L/min}$ and
$195 \times 0.160 \times 0.150 \approx \qty{4.7}{L/min}$. The stroke
volume is what training enlarges: the second subject is the endurance
trained one.
\end{solution}

\begin{solution}{exo:g10:heart-lungs-effort:13}
Oxygen breathed in: $0.21 \times 100 = \qty{21}{L/min}$; used:
\qty{3.5}{L/min}, i.e. 17\%. The rest is breathed out again --- expired
air still holds about 16\% oxygen, which is why mouth-to-mouth
resuscitation works.
\end{solution}

\begin{solution}{exo:g10:heart-lungs-effort:14}
Output at most about $110 \times 0.12 \approx \qty{13}{L/min}$ instead
of 24: the $\dot V\!\mathrm{O_2}$max falls by nearly half, to about
\qty{2}{L/min}. The patient can walk and cycle gently but not sustain
any hard effort.
\end{solution}

\begin{solution}{exo:g10:heart-lungs-effort:15}
Delivery per litre down by a quarter, extraction unchanged in
proportion: $\dot V\!\mathrm{O_2}$max falls by 25\%. Over the
following weeks the body makes more red blood \omterm{def:g10:cells-common-unit:cell}{cells}, raising the
oxygen carried per litre back towards \qty{200}{mL}.
\end{solution}

\begin{solution}{pb:g10:heart-lungs-effort:1}
\textbf{1.} $62 \times 0.072 \approx 4.5$; $105 \times 0.100 = 10.5$;
$140 \times 0.116 \approx 16.2$; $175 \times 0.120 = 21.0$;
$198 \times 0.120 \approx \qty{23.8}{L/min}$.

\textbf{2.} $23.8/4.5 \approx 5.3$: the rate contributes $198/62
\approx 3.2$, the stroke volume $120/72 \approx 1.7$, and $3.2 \times
1.7 \approx 5.3$.

\textbf{3.} Above about \qty{225}{W} the stroke volume stays at
\qty{120}{mL}; only the heart rate raises the output further.

\textbf{4.} $220 - 17 = 203$: his 198 is within the rule's precision.

\textbf{5.} $5/4.5 \approx \qty{1.1}{min}$, about \qty{67}{s}; at
\qty{300}{W}, $5/23.8 \approx \qty{0.21}{min}$, about \qty{13}{s}.

\textbf{6.} $0.5 \times 13 = 6.5$; $1.2 \times 18 = 21.6$; $1.8 \times
24 = 43.2$; $2.3 \times 32 = 73.6$; $2.6 \times 44 \approx
\qty{114}{L/min}$.

\textbf{7.} $114/6.5 \approx 18$: the tidal volume rose 5.2 times, the
rate 3.4 times; the volume contributes more.

\textbf{8.} $2.6/5.0 = 52\%$ of the vital capacity per breath.

\textbf{9.} No: \omterm{def:g10:heart-lungs-effort:ventilation}{ventilation} has a large reserve and the blood leaving
the lungs stays saturated. The limit is the heart's stroke volume,
which sets how much oxygen each beat delivers.

\textbf{10.} $0.7 \times 6.5 \approx \qty{4.6}{L/min}$ and $0.7 \times
114 \approx \qty{80}{L/min}$.

\textbf{11.} $4.5 \times 0.052 \approx 0.23$; $10.5 \times 0.090
\approx 0.95$; $16.2 \times 0.120 \approx 1.94$; $21.0 \times 0.145
\approx 3.05$; $23.8 \times 0.155 \approx \qty{3.7}{L/min}$.

\textbf{12.} $3700/68 \approx \qty{54}{mL/min}$ per kilogram. The rate
has reached its maximum and the increase from 225 to \qty{300}{W} is
small: this is his $\dot V\!\mathrm{O_2}$max, that of a fit, trained
student.

\textbf{13.} $3.7/0.23 \approx 16$: the heart's factor 5.3 multiplied
by the extraction's $155/52 \approx 3.0$.

\textbf{14.} $3.7 \times 20 = \qty{74}{kJ}$ per minute, about
\qty{1230}{W}; at rest $0.23 \times 20/60 \approx \qty{77}{W}$; above
rest about \qty{1150}{W}, so the efficiency is $300/1150 \approx
26\%$.

\textbf{15.} Roughly a straight line from \num{0.23} to
\qty{3.05}{L/min} between 0 and \qty{225}{W}, then bending: the last
\qty{75}{W} add only \qty{0.65}{L/min}. The ceiling is about
\qty{3.7}{L/min}.

\textbf{16.} Anticipation: the brain's expectation of the effort is
relayed to the cardiac centres and the adrenal glands, raising the
rate before any movement.

\textbf{17.} The \omterm{def:g10:heart-lungs-effort:ventilation}{alveoli} oxygenate the blood fully even at maximal
\omterm{def:g10:heart-lungs-effort:ventilation}{ventilation} and output; the lungs are not the limit. The limit is the
amount of blood the heart can send per minute.

\textbf{18.} $20/23.8 \approx 84\%$. The gut and kidneys, which took
half of the resting output, now receive about \qty{1}{L/min}; their
arteries have narrowed and the flow has been redirected to the
muscles.

\textbf{19.} $198 \times 0.140 \times 0.155 \approx \qty{4.3}{L/min}$,
about \qty{63}{mL/min} per kilogram.

\textbf{20.} \omterm{def:g10:heart-lungs-effort:output}{Cardiac output} multiplied by 5.3, oxygen extraction by 3,
\omterm{def:g10:heart-lungs-effort:ventilation}{ventilation} by 18; training changed the heart's stroke volume, and
with it the ceiling.
\end{solution}
